\documentclass[reqno]{amsart} 
\usepackage{amssymb,latexsym,amsmath,amscd,graphicx,setspace,amsthm,verbatim,enumitem}
\usepackage[margin = 3 cm]{geometry}

\def\upint{\mathchoice%
    {\mkern13mu\overline{\vphantom{\intop}\mkern7mu}\mkern-20mu}%
    {\mkern7mu\overline{\vphantom{\intop}\mkern7mu}\mkern-14mu}%
    {\mkern7mu\overline{\vphantom{\intop}\mkern7mu}\mkern-14mu}%
    {\mkern7mu\overline{\vphantom{\intop}\mkern7mu}\mkern-14mu}%
  \int}
\def\lowint{\mkern3mu\underline{\vphantom{\intop}\mkern7mu}\mkern-10mu\int}


\theoremstyle{plain}
\newtheorem{theorem}{Theorem}[section]
\newtheorem{proposition}{Proposition}
\newtheorem{corollary}{Corollary}
\newtheorem{lemma}{Lemma}
\newtheorem{conjecture}{Conjecture}
\newtheorem{question}{Question}
\newtheorem{problem}{Problem}
      
\theoremstyle{definition}
\newtheorem{definition}{Definition}

\newenvironment{solution1}{\paragraph{\emph{Solution $1$}.}}{\hfill$\square$}
\newenvironment{solution2}{\paragraph{\emph{Solution $2$}.}}{\hfill$\square$}
\newenvironment{solution3}{\paragraph{\emph{Solution $3$}.}}{\hfill$\square$}

\begin{document} 

\title[Homework 11]{Homework 11}

\date{\today} 
\maketitle 






\begin{problem}
In class, we went over the following example of a pointwise converging sequence of functions.  The sequence was $f_{n}:[0,1] \rightarrow [0,\infty]$ defined via
\begin{equation*}
x \mapsto f_{n}(x) = 
\begin{cases}
n,& \text{ if } 0 \le x \le \frac{1}{n} ; \\
1-x, & \text{ if } \frac{1}{n} \le x \le 1.
\end{cases}
\end{equation*}
\begin{enumerate}
\item Are the functions $f_{n}$ measurable?  Explain.
\item Does the sequence $(f_{n})_{n=1}^{\infty}$ converge pointwise to a function $f:[0,1] \rightarrow [0,\infty]$?  If no, explain, if yes, find $f$.
\item Is the function $f$ measurable?  Explain.
\item Calculate $\int_{[0,1]}f_{n}$.
\item Calculate $\int_{[0,1]}f$.
\item Does one have
$$\int_{[0,1]}f_{n} \rightarrow \int_{[0,1]}f, $$
when $n \to \infty$?
\item Is this a counterexample to the monotone convergence theorem?  Explain.
\end{enumerate}
\end{problem}
\begin{solution1}

\end{solution1}


\begin{problem}
Let $(X,\mathcal{A})$ be a measurable space and let $f:X \rightarrow [-\infty,\infty]$ be a function.  In class, we will define two functions $f^{+},f^{-}:X \rightarrow [0,\infty]$ via
\begin{equation*}
x \mapsto f^{+}(x) = 
\begin{cases}
f(x), & \text{ if } f(x) \ge 0;\\
0, & \text{ if } f(x) < 0.
\end{cases}
\end{equation*}
and 
\begin{equation*}
x \mapsto f^{-}(x) = 
\begin{cases}
0, & \text{ if } f(x) \ge 0;\\
-f(x), & \text{ if } f(x) < 0.
\end{cases}
\end{equation*}
Show that $f$ is measurable if and only if both $f^{+}$ and $f^{-}$ are measurable.
\end{problem}

\begin{solution1}

\end{solution1}


We proved a similar result to the following problem in class under the assumption $\mu(X)< \infty$ and $f: X \rightarrow \mathbb{R}$ is bounded.
\begin{problem}[3A-3]
Suppose $(X,\mathcal{A},\mu)$ is a measure space and $f:X \rightarrow [0,\infty]$ is a measurable function.  Prove that
$$\int_{X}f > 0 \text{ if and only if } \mu(\{x \in X : f(x) > 0\})> 0. $$
\end{problem}
\begin{solution1}

\end{solution1}

\begin{problem}[3A-9]
Let $(X,\mathcal{A},\mu)$ be a measure space and let $f: X \rightarrow [0,\infty]$ be a measurable function.  Define $\nu:\mathcal{A} \rightarrow [0,\infty]$ by
$$\nu(A) = \int_{X} f \cdot \chi_{A} $$
for $A \in \mathcal{A}$.  Prove that $\nu$ is a measure on $(X,\mathcal{A})$.
\end{problem}
\begin{solution1}

\end{solution1}





\begin{problem}[3A-10]
Suppose $(X,\mathcal{A},\mu)$ is a measure space and let $f_{n}: X \rightarrow [0,\infty]$ be a sequence of measurable functions.  Define $f:X \rightarrow [0,\infty]$ by
$$f(x) = \sum_{k=1}^{\infty}f_{k}(x). $$
Prove that
$$\int_{X} f = \sum_{k=1}^{\infty} \int_{X}f_{k}. $$
\end{problem}
\begin{solution1}

\end{solution1}





\end{document}









